% --- Archon physics formalization source begin ---
% archon:physics
% archon:covers PhyXMiniProblems/problem_phyx_mini_0895.lean
% archon:source-report reports/phyx_mini/problem_phyx_mini_0895.source.json

\chapter{Physics problem phyx\_mini\_0895}
\label{ch:PhyXMiniProblems_problem_phyx_mini_0895}

\paragraph{Problem source.}
\#\# Question

What is the magnitude of the force \textbackslash{}( \textbackslash{}vec\{F\} \textbackslash{}) on the \textbackslash{}( -10\textbackslash{},\textbackslash{}mathrm\{nC\} \textbackslash{}) charge in figure?

\#\# Answer choices

- A: A: 1.25 \textbackslash{}times 10\textasciicircum{}\{-3\}N

- B: B: 1.35 \textbackslash{}times 10\textasciicircum{}\{-3\}N

- C: C: 4.3 \textbackslash{}times 10\textasciicircum{}\{-3\}N

- D: D: 1.33 \textbackslash{}times 10\textasciicircum{}\{-3\}N

\#\# Figure caption (auxiliary; use the image as primary evidence)

The image illustrates three charged particles arranged in a right-angle triangle. 

- The top left particle is a red circle with a positive charge symbol "+" labeled as "15 nC."
- The top right particle is a teal circle with a negative charge symbol "−" labeled as "−5.0 nC."
- The bottom right particle is another teal circle with a negative charge symbol "−" labeled as "−10 nC."

The distances between the particles are indicated by dashed lines:
- The horizontal distance between the red and top teal circles is labeled "3.0 cm."
- The vertical distance between the two teal circles is labeled "1.0 cm."

The layout forms a right-angle triangle with the 90-degree angle at the top right corner.

\#\# Dataset metadata

Electrostatics; Physical Model Grounding Reasoning, Multi-Formula Reasoning

\paragraph{Recorded answer/context.}
C: C: 4.3 \textbackslash{}times 10\textasciicircum{}\{-3\}N

\paragraph{Figure/image path.}
/root/proposal\_for\_physic/science-mango/phyx\_mini\_run/phyx\_data/test\_image/895.png

\paragraph{Formalization target.}
create a compiling Lean file with sorry bodies at `PhyXMiniProblems/problem\_phyx\_mini\_0895.lean`. The Lean declarations must preserve the physical quantities, dimensions or dimensional roles, figure labels, governing-law hypotheses, and final relation expressed by this problem.
Use Mathlib/Physlib names found through LeanExplore where available. If a physics API is missing, introduce faithful local abstractions rather than scalar placeholder aliases.

\begin{theorem}[Physics formalization target]
\label{thm:physics:phyx_mini_0895:target}
\lean{PhyXMiniProblems.ProblemPhyXMini0895.problem_phyx_mini_0895}
\uses{def:physics:phyx-mini-0895:phyxminiproblems-problemphyxmini0895-forcemagnitudeinnewtons, def:physics:phyx-mini-0895:phyxminiproblems-problemphyxmini0895-threepointchargesetup, def:physics:phyx-mini-0895:phyxminiproblems-problemphyxmini0895-matchessuppliedthreechargefigure, def:physics:phyx-mini-0895:phyxminiproblems-problemphyxmini0895-usesschoolcoulombconstant, def:physics:phyx-mini-0895:phyxminiproblems-problemphyxmini0895-hasphysicalthreechargeparameters, def:physics:phyx-mini-0895:phyxminiproblems-problemphyxmini0895-satisfiesvectorcoulombforcelaw, def:physics:phyx-mini-0895:phyxminiproblems-problemphyxmini0895-satisfieselectrostaticforcesuperposition, def:physics:phyx-mini-0895:phyxminiproblems-problemphyxmini0895-satisfiesresultantforcemagnitudelaw, def:physics:phyx-mini-0895:phyxminiproblems-problemphyxmini0895-displayedforcemagnitudeinnewtons, def:physics:phyx-mini-0895:phyxminiproblems-problemphyxmini0895-roundstoatresolution, def:physics:phyx-mini-0895:phyxminiproblems-problemphyxmini0895-isuniqueclosestdisplayedforcechoice}
The assigned autoformalize agent should translate this physics problem into Lean declarations in the covered file, with theorem and lemma proof bodies written as `by sorry`.
\end{theorem}
\begin{proof}
This is an autoformalization task, not a proof task. Produce statements that can later be proved by the physics prover without weakening the physical model.
\end{proof}

% --- Archon declaration topology begin ---
\section{Declaration topology}
\begin{definition}[force Dimension]
  \label{def:physics:phyx-mini-0895:phyxminiproblems-problemphyxmini0895-forcedimension}
  \lean{PhyXMiniProblems.ProblemPhyXMini0895.forceDimension}
  The physical dimension M L T⁻² of force
\end{definition}
\begin{proof}
  This is the declared model definition of force Dimension.
\end{proof}
\begin{definition}[Length Quantity]
  \label{def:physics:phyx-mini-0895:phyxminiproblems-problemphyxmini0895-lengthquantity}
  \lean{PhyXMiniProblems.ProblemPhyXMini0895.LengthQuantity}
  A nonnegative, unit-independent physical length
\end{definition}
\begin{proof}
  This is the declared model definition of Length Quantity.
\end{proof}
\begin{definition}[Signed Charge Quantity]
  \label{def:physics:phyx-mini-0895:phyxminiproblems-problemphyxmini0895-signedchargequantity}
  \lean{PhyXMiniProblems.ProblemPhyXMini0895.SignedChargeQuantity}
  A signed, unit-independent electric charge
\end{definition}
\begin{proof}
  This is the declared model definition of Signed Charge Quantity.
\end{proof}
\begin{definition}[Planar Position Quantity]
  \label{def:physics:phyx-mini-0895:phyxminiproblems-problemphyxmini0895-planarpositionquantity}
  \lean{PhyXMiniProblems.ProblemPhyXMini0895.PlanarPositionQuantity}
  A unit-independent planar position
\end{definition}
\begin{proof}
  This is the declared model definition of Planar Position Quantity.
\end{proof}
\begin{definition}[Planar Force Quantity]
  \label{def:physics:phyx-mini-0895:phyxminiproblems-problemphyxmini0895-planarforcequantity}
  \lean{PhyXMiniProblems.ProblemPhyXMini0895.PlanarForceQuantity}
  \uses{def:physics:phyx-mini-0895:phyxminiproblems-problemphyxmini0895-forcedimension}
  A unit-independent planar force vector
\end{definition}
\begin{proof}
  This is the declared model definition of Planar Force Quantity.
\end{proof}
\begin{definition}[Force Magnitude Quantity]
  \label{def:physics:phyx-mini-0895:phyxminiproblems-problemphyxmini0895-forcemagnitudequantity}
  \lean{PhyXMiniProblems.ProblemPhyXMini0895.ForceMagnitudeQuantity}
  \uses{def:physics:phyx-mini-0895:phyxminiproblems-problemphyxmini0895-forcedimension}
  A nonnegative, unit-independent force magnitude
\end{definition}
\begin{proof}
  This is the declared model definition of Force Magnitude Quantity.
\end{proof}
\begin{definition}[length In Meters]
  \label{def:physics:phyx-mini-0895:phyxminiproblems-problemphyxmini0895-lengthinmeters}
  \lean{PhyXMiniProblems.ProblemPhyXMini0895.lengthInMeters}
  \uses{def:physics:phyx-mini-0895:phyxminiproblems-problemphyxmini0895-lengthquantity}
  Coherent-SI readout of a length in metres
\end{definition}
\begin{proof}
  This is the declared model definition of length In Meters.
\end{proof}
\begin{definition}[length In Centimeters]
  \label{def:physics:phyx-mini-0895:phyxminiproblems-problemphyxmini0895-lengthincentimeters}
  \lean{PhyXMiniProblems.ProblemPhyXMini0895.lengthInCentimeters}
  \uses{def:physics:phyx-mini-0895:phyxminiproblems-problemphyxmini0895-lengthquantity, def:physics:phyx-mini-0895:phyxminiproblems-problemphyxmini0895-lengthinmeters}
  Centimetre readout used by the two dimension labels in the image
\end{definition}
\begin{proof}
  This is the declared model definition of length In Centimeters.
\end{proof}
\begin{definition}[charge In Coulombs]
  \label{def:physics:phyx-mini-0895:phyxminiproblems-problemphyxmini0895-chargeincoulombs}
  \lean{PhyXMiniProblems.ProblemPhyXMini0895.chargeInCoulombs}
  \uses{def:physics:phyx-mini-0895:phyxminiproblems-problemphyxmini0895-signedchargequantity}
  Coherent-SI readout of a signed charge in coulombs
\end{definition}
\begin{proof}
  This is the declared model definition of charge In Coulombs.
\end{proof}
\begin{definition}[charge In Nanocoulombs]
  \label{def:physics:phyx-mini-0895:phyxminiproblems-problemphyxmini0895-chargeinnanocoulombs}
  \lean{PhyXMiniProblems.ProblemPhyXMini0895.chargeInNanocoulombs}
  \uses{def:physics:phyx-mini-0895:phyxminiproblems-problemphyxmini0895-signedchargequantity, def:physics:phyx-mini-0895:phyxminiproblems-problemphyxmini0895-chargeincoulombs}
  Nanocoulomb readout used by the three printed charge labels
\end{definition}
\begin{proof}
  This is the declared model definition of charge In Nanocoulombs.
\end{proof}
\begin{definition}[position Vector In Meters]
  \label{def:physics:phyx-mini-0895:phyxminiproblems-problemphyxmini0895-positionvectorinmeters}
  \lean{PhyXMiniProblems.ProblemPhyXMini0895.positionVectorInMeters}
  \uses{def:physics:phyx-mini-0895:phyxminiproblems-problemphyxmini0895-planarpositionquantity}
  Coherent-SI Cartesian position vector, in metres
\end{definition}
\begin{proof}
  This is the declared model definition of position Vector In Meters.
\end{proof}
\begin{definition}[force Vector In Newtons]
  \label{def:physics:phyx-mini-0895:phyxminiproblems-problemphyxmini0895-forcevectorinnewtons}
  \lean{PhyXMiniProblems.ProblemPhyXMini0895.forceVectorInNewtons}
  \uses{def:physics:phyx-mini-0895:phyxminiproblems-problemphyxmini0895-planarforcequantity}
  Coherent-SI Cartesian force vector, in newtons
\end{definition}
\begin{proof}
  This is the declared model definition of force Vector In Newtons.
\end{proof}
\begin{definition}[force Magnitude In Newtons]
  \label{def:physics:phyx-mini-0895:phyxminiproblems-problemphyxmini0895-forcemagnitudeinnewtons}
  \lean{PhyXMiniProblems.ProblemPhyXMini0895.forceMagnitudeInNewtons}
  \uses{def:physics:phyx-mini-0895:phyxminiproblems-problemphyxmini0895-forcemagnitudequantity}
  Coherent-SI readout of a physical force magnitude, in newtons
\end{definition}
\begin{proof}
  This is the declared model definition of force Magnitude In Newtons.
\end{proof}
\begin{definition}[Charge Site]
  \label{def:physics:phyx-mini-0895:phyxminiproblems-problemphyxmini0895-chargesite}
  \lean{PhyXMiniProblems.ProblemPhyXMini0895.ChargeSite}
  The three corners occupied by charges in the primary image
\end{definition}
\begin{proof}
  This is the declared model definition of Charge Site.
\end{proof}
\begin{definition}[Particle]
  \label{def:physics:phyx-mini-0895:phyxminiproblems-problemphyxmini0895-particle}
  \lean{PhyXMiniProblems.ProblemPhyXMini0895.Particle}
  The three particles, named by physical role and image location
\end{definition}
\begin{proof}
  This is the declared model definition of Particle.
\end{proof}
\begin{definition}[Force Source]
  \label{def:physics:phyx-mini-0895:phyxminiproblems-problemphyxmini0895-forcesource}
  \lean{PhyXMiniProblems.ProblemPhyXMini0895.ForceSource}
  The two particles which exert force on the target particle
\end{definition}
\begin{proof}
  This is the declared model definition of Force Source.
\end{proof}
\begin{definition}[source Particle]
  \label{def:physics:phyx-mini-0895:phyxminiproblems-problemphyxmini0895-sourceparticle}
  \lean{PhyXMiniProblems.ProblemPhyXMini0895.sourceParticle}
  \uses{def:physics:phyx-mini-0895:phyxminiproblems-problemphyxmini0895-particle, def:physics:phyx-mini-0895:phyxminiproblems-problemphyxmini0895-forcesource}
  Particle associated with each source-force contribution
\end{definition}
\begin{proof}
  This is the declared model definition of source Particle.
\end{proof}
\begin{definition}[expected Particle Site]
  \label{def:physics:phyx-mini-0895:phyxminiproblems-problemphyxmini0895-expectedparticlesite}
  \lean{PhyXMiniProblems.ProblemPhyXMini0895.expectedParticleSite}
  \uses{def:physics:phyx-mini-0895:phyxminiproblems-problemphyxmini0895-chargesite, def:physics:phyx-mini-0895:phyxminiproblems-problemphyxmini0895-particle}
  Site occupied by each named particle in the supplied raster
\end{definition}
\begin{proof}
  This is the declared model definition of expected Particle Site.
\end{proof}
\begin{definition}[expected Charge In Nanocoulombs]
  \label{def:physics:phyx-mini-0895:phyxminiproblems-problemphyxmini0895-expectedchargeinnanocoulombs}
  \lean{PhyXMiniProblems.ProblemPhyXMini0895.expectedChargeInNanocoulombs}
  \uses{def:physics:phyx-mini-0895:phyxminiproblems-problemphyxmini0895-particle}
  Signed nanocoulomb value printed beside each particle
\end{definition}
\begin{proof}
  This is the declared model definition of expected Charge In Nanocoulombs.
\end{proof}
\begin{definition}[Figure Charge Color]
  \label{def:physics:phyx-mini-0895:phyxminiproblems-problemphyxmini0895-figurechargecolor}
  \lean{PhyXMiniProblems.ProblemPhyXMini0895.FigureChargeColor}
  Colour used for a charge circle in the supplied figure
\end{definition}
\begin{proof}
  This is the declared model definition of Figure Charge Color.
\end{proof}
\begin{definition}[Figure Charge Sign]
  \label{def:physics:phyx-mini-0895:phyxminiproblems-problemphyxmini0895-figurechargesign}
  \lean{PhyXMiniProblems.ProblemPhyXMini0895.FigureChargeSign}
  Sign glyph drawn inside a charge circle
\end{definition}
\begin{proof}
  This is the declared model definition of Figure Charge Sign.
\end{proof}
\begin{definition}[expected Particle Color]
  \label{def:physics:phyx-mini-0895:phyxminiproblems-problemphyxmini0895-expectedparticlecolor}
  \lean{PhyXMiniProblems.ProblemPhyXMini0895.expectedParticleColor}
  \uses{def:physics:phyx-mini-0895:phyxminiproblems-problemphyxmini0895-particle, def:physics:phyx-mini-0895:phyxminiproblems-problemphyxmini0895-figurechargecolor}
  Expected circle colour of each particle
\end{definition}
\begin{proof}
  This is the declared model definition of expected Particle Color.
\end{proof}
\begin{definition}[expected Particle Sign]
  \label{def:physics:phyx-mini-0895:phyxminiproblems-problemphyxmini0895-expectedparticlesign}
  \lean{PhyXMiniProblems.ProblemPhyXMini0895.expectedParticleSign}
  \uses{def:physics:phyx-mini-0895:phyxminiproblems-problemphyxmini0895-particle, def:physics:phyx-mini-0895:phyxminiproblems-problemphyxmini0895-figurechargesign}
  Expected sign glyph of each particle
\end{definition}
\begin{proof}
  This is the declared model definition of expected Particle Sign.
\end{proof}
\begin{definition}[Labelled Separation]
  \label{def:physics:phyx-mini-0895:phyxminiproblems-problemphyxmini0895-labelledseparation}
  \lean{PhyXMiniProblems.ProblemPhyXMini0895.LabelledSeparation}
  The two dashed charge-to-charge segments carrying numerical labels
\end{definition}
\begin{proof}
  This is the declared model definition of Labelled Separation.
\end{proof}
\begin{definition}[Three Charge Figure]
  \label{def:physics:phyx-mini-0895:phyxminiproblems-problemphyxmini0895-threechargefigure}
  \lean{PhyXMiniProblems.ProblemPhyXMini0895.ThreeChargeFigure}
  \uses{def:physics:phyx-mini-0895:phyxminiproblems-problemphyxmini0895-chargesite, def:physics:phyx-mini-0895:phyxminiproblems-problemphyxmini0895-particle, def:physics:phyx-mini-0895:phyxminiproblems-problemphyxmini0895-figurechargecolor, def:physics:phyx-mini-0895:phyxminiproblems-problemphyxmini0895-figurechargesign, def:physics:phyx-mini-0895:phyxminiproblems-problemphyxmini0895-labelledseparation}
  Literal presentation data transcribed from image 895.png
\end{definition}
\begin{proof}
  This is the declared model definition of Three Charge Figure.
\end{proof}
\begin{definition}[Three Point Charge Setup]
  \label{def:physics:phyx-mini-0895:phyxminiproblems-problemphyxmini0895-threepointchargesetup}
  \lean{PhyXMiniProblems.ProblemPhyXMini0895.ThreePointChargeSetup}
  \uses{def:physics:phyx-mini-0895:phyxminiproblems-problemphyxmini0895-lengthquantity, def:physics:phyx-mini-0895:phyxminiproblems-problemphyxmini0895-signedchargequantity, def:physics:phyx-mini-0895:phyxminiproblems-problemphyxmini0895-planarpositionquantity, def:physics:phyx-mini-0895:phyxminiproblems-problemphyxmini0895-planarforcequantity, def:physics:phyx-mini-0895:phyxminiproblems-problemphyxmini0895-forcemagnitudequantity, def:physics:phyx-mini-0895:phyxminiproblems-problemphyxmini0895-particle, def:physics:phyx-mini-0895:phyxminiproblems-problemphyxmini0895-forcesource, def:physics:phyx-mini-0895:phyxminiproblems-problemphyxmini0895-threechargefigure}
  Independent physical quantities and force observables in the problem. In particular, the pairwise forces, resultant force, and resultant magnitude are not defined from an answer choice. They are constrained only by the governing-law predicates below
\end{definition}
\begin{proof}
  This is the declared model definition of Three Point Charge Setup.
\end{proof}
\begin{definition}[displacement From Source To Target In Meters]
  \label{def:physics:phyx-mini-0895:phyxminiproblems-problemphyxmini0895-displacementfromsourcetotargetinmeters}
  \lean{PhyXMiniProblems.ProblemPhyXMini0895.displacementFromSourceToTargetInMeters}
  \uses{def:physics:phyx-mini-0895:phyxminiproblems-problemphyxmini0895-positionvectorinmeters, def:physics:phyx-mini-0895:phyxminiproblems-problemphyxmini0895-forcesource, def:physics:phyx-mini-0895:phyxminiproblems-problemphyxmini0895-sourceparticle, def:physics:phyx-mini-0895:phyxminiproblems-problemphyxmini0895-threepointchargesetup}
  Displacement from a source particle to the target particle, in metres
\end{definition}
\begin{proof}
  This is the declared model definition of displacement From Source To Target In Meters.
\end{proof}
\begin{definition}[Matches Supplied Three Charge Figure]
  \label{def:physics:phyx-mini-0895:phyxminiproblems-problemphyxmini0895-matchessuppliedthreechargefigure}
  \lean{PhyXMiniProblems.ProblemPhyXMini0895.MatchesSuppliedThreeChargeFigure}
  \uses{def:physics:phyx-mini-0895:phyxminiproblems-problemphyxmini0895-lengthinmeters, def:physics:phyx-mini-0895:phyxminiproblems-problemphyxmini0895-lengthincentimeters, def:physics:phyx-mini-0895:phyxminiproblems-problemphyxmini0895-chargeinnanocoulombs, def:physics:phyx-mini-0895:phyxminiproblems-problemphyxmini0895-positionvectorinmeters, def:physics:phyx-mini-0895:phyxminiproblems-problemphyxmini0895-expectedparticlesite, def:physics:phyx-mini-0895:phyxminiproblems-problemphyxmini0895-expectedchargeinnanocoulombs, def:physics:phyx-mini-0895:phyxminiproblems-problemphyxmini0895-expectedparticlecolor, def:physics:phyx-mini-0895:phyxminiproblems-problemphyxmini0895-expectedparticlesign, def:physics:phyx-mini-0895:phyxminiproblems-problemphyxmini0895-threepointchargesetup}
  All unambiguous qualitative and numerical data transcribed from the primary raster, together with its association to the independent physical charges, lengths, and positions. No force magnitude appears here
\end{definition}
\begin{proof}
  This is the declared model definition of Matches Supplied Three Charge Figure.
\end{proof}
\begin{definition}[Uses School Coulomb Constant]
  \label{def:physics:phyx-mini-0895:phyxminiproblems-problemphyxmini0895-usesschoolcoulombconstant}
  \lean{PhyXMiniProblems.ProblemPhyXMini0895.UsesSchoolCoulombConstant}
  \uses{def:physics:phyx-mini-0895:phyxminiproblems-problemphyxmini0895-threepointchargesetup}
  The rounded Coulomb constant used in the standard multiple-choice calculation. Physlib's Electromagnetism.EMSystem.coulombConstant is the coherent-SI scalar 1 / (4 π ε₀) appearing in the vector force law
\end{definition}
\begin{proof}
  This is the declared model definition of Uses School Coulomb Constant.
\end{proof}
\begin{definition}[Has Physical Three Charge Parameters]
  \label{def:physics:phyx-mini-0895:phyxminiproblems-problemphyxmini0895-hasphysicalthreechargeparameters}
  \lean{PhyXMiniProblems.ProblemPhyXMini0895.HasPhysicalThreeChargeParameters}
  \uses{def:physics:phyx-mini-0895:phyxminiproblems-problemphyxmini0895-lengthinmeters, def:physics:phyx-mini-0895:phyxminiproblems-problemphyxmini0895-threepointchargesetup, def:physics:phyx-mini-0895:phyxminiproblems-problemphyxmini0895-displacementfromsourcetotargetinmeters}
  Positivity and source-separation conditions selecting the physical case
\end{definition}
\begin{proof}
  This is the declared model definition of Has Physical Three Charge Parameters.
\end{proof}
\begin{definition}[Satisfies Vector Coulomb Force Law]
  \label{def:physics:phyx-mini-0895:phyxminiproblems-problemphyxmini0895-satisfiesvectorcoulombforcelaw}
  \lean{PhyXMiniProblems.ProblemPhyXMini0895.SatisfiesVectorCoulombForceLaw}
  \uses{def:physics:phyx-mini-0895:phyxminiproblems-problemphyxmini0895-chargeincoulombs, def:physics:phyx-mini-0895:phyxminiproblems-problemphyxmini0895-forcevectorinnewtons, def:physics:phyx-mini-0895:phyxminiproblems-problemphyxmini0895-sourceparticle, def:physics:phyx-mini-0895:phyxminiproblems-problemphyxmini0895-threepointchargesetup, def:physics:phyx-mini-0895:phyxminiproblems-problemphyxmini0895-displacementfromsourcetotargetinmeters}
  For source charge qₛ, target charge qₜ, and displacement r from source to target, the force on the target is F = k qₜ qₛ r / ‖r‖³. The signed product of charges supplies attraction or repulsion. This general pairwise law contains no requested numerical resultant
\end{definition}
\begin{proof}
  This is the declared model definition of Satisfies Vector Coulomb Force Law.
\end{proof}
\begin{definition}[Satisfies Electrostatic Force Superposition]
  \label{def:physics:phyx-mini-0895:phyxminiproblems-problemphyxmini0895-satisfieselectrostaticforcesuperposition}
  \lean{PhyXMiniProblems.ProblemPhyXMini0895.SatisfiesElectrostaticForceSuperposition}
  \uses{def:physics:phyx-mini-0895:phyxminiproblems-problemphyxmini0895-forcevectorinnewtons, def:physics:phyx-mini-0895:phyxminiproblems-problemphyxmini0895-forcesource, def:physics:phyx-mini-0895:phyxminiproblems-problemphyxmini0895-threepointchargesetup}
  The resultant force is the vector sum of the two pairwise forces
\end{definition}
\begin{proof}
  This is the declared model definition of Satisfies Electrostatic Force Superposition.
\end{proof}
\begin{definition}[Satisfies Resultant Force Magnitude Law]
  \label{def:physics:phyx-mini-0895:phyxminiproblems-problemphyxmini0895-satisfiesresultantforcemagnitudelaw}
  \lean{PhyXMiniProblems.ProblemPhyXMini0895.SatisfiesResultantForceMagnitudeLaw}
  \uses{def:physics:phyx-mini-0895:phyxminiproblems-problemphyxmini0895-forcevectorinnewtons, def:physics:phyx-mini-0895:phyxminiproblems-problemphyxmini0895-forcemagnitudeinnewtons, def:physics:phyx-mini-0895:phyxminiproblems-problemphyxmini0895-threepointchargesetup}
  The independent nonnegative magnitude observable is the Euclidean norm of the resultant planar force vector
\end{definition}
\begin{proof}
  This is the declared model definition of Satisfies Resultant Force Magnitude Law.
\end{proof}
\begin{definition}[Answer Choice]
  \label{def:physics:phyx-mini-0895:phyxminiproblems-problemphyxmini0895-answerchoice}
  \lean{PhyXMiniProblems.ProblemPhyXMini0895.AnswerChoice}
  Labels attached to the four force-magnitude choices in the source
\end{definition}
\begin{proof}
  This is the declared model definition of Answer Choice.
\end{proof}
\begin{definition}[displayed Force Magnitude In Newtons]
  \label{def:physics:phyx-mini-0895:phyxminiproblems-problemphyxmini0895-displayedforcemagnitudeinnewtons}
  \lean{PhyXMiniProblems.ProblemPhyXMini0895.displayedForceMagnitudeInNewtons}
  \uses{def:physics:phyx-mini-0895:phyxminiproblems-problemphyxmini0895-answerchoice}
  Printed magnitude beside each answer choice, in newtons
\end{definition}
\begin{proof}
  This is the declared model definition of displayed Force Magnitude In Newtons.
\end{proof}
\begin{definition}[Rounds To At Resolution]
  \label{def:physics:phyx-mini-0895:phyxminiproblems-problemphyxmini0895-roundstoatresolution}
  \lean{PhyXMiniProblems.ProblemPhyXMini0895.RoundsToAtResolution}
  value rounds to displayed at a stated display resolution when it lies strictly within half a resolution step. For 4.3 * 10⁻³ N, the last printed digit has resolution 0.1 * 10⁻³ N
\end{definition}
\begin{proof}
  This is the declared model definition of Rounds To At Resolution.
\end{proof}
\begin{definition}[Is Unique Closest Displayed Force Choice]
  \label{def:physics:phyx-mini-0895:phyxminiproblems-problemphyxmini0895-isuniqueclosestdisplayedforcechoice}
  \lean{PhyXMiniProblems.ProblemPhyXMini0895.IsUniqueClosestDisplayedForceChoice}
  \uses{def:physics:phyx-mini-0895:phyxminiproblems-problemphyxmini0895-forcemagnitudeinnewtons, def:physics:phyx-mini-0895:phyxminiproblems-problemphyxmini0895-threepointchargesetup, def:physics:phyx-mini-0895:phyxminiproblems-problemphyxmini0895-answerchoice, def:physics:phyx-mini-0895:phyxminiproblems-problemphyxmini0895-displayedforcemagnitudeinnewtons}
  A choice is uniquely closest to the independently modelled magnitude
\end{definition}
\begin{proof}
  This is the declared model definition of Is Unique Closest Displayed Force Choice.
\end{proof}
% --- Archon declaration topology end ---
% --- Archon physics formalization source end ---
